% Verbatim copy of /Users/alingalatan/workspaces/grothe/handoff/research-degree5/ai-notes/chatgpt-thread2-gaussian_domination_note.tex (file dated 2026-11-05). Only this line was added.
% output/pdf/gaussian_domination_note.tex
\documentclass[11pt]{article}
\usepackage[T1]{fontenc}
\usepackage{lmodern,microtype,amsmath,amssymb,mathtools}
\usepackage[a4paper,top=23mm,bottom=23mm,left=25mm,right=25mm,headheight=14pt]{geometry}
\usepackage[table]{xcolor}
\usepackage{fancyhdr}
\usepackage[colorlinks=true,urlcolor=ink,linkcolor=ink]{hyperref}
\definecolor{ink}{HTML}{233B53}
\definecolor{muted}{HTML}{505C65}
\pagestyle{fancy}
\fancyhf{}
\fancyhead[L]{\small\color{muted} Gaussian domination in cubic chaos}
\fancyhead[R]{\small\color{muted} 5 November 2026}
\fancyfoot[C]{\small\thepage}
\renewcommand{\headrulewidth}{0.3pt}
\setlength{\parindent}{0pt}
\setlength{\parskip}{4pt}
\setlength{\abovedisplayskip}{5pt plus 2pt minus 1pt}
\setlength{\belowdisplayskip}{5pt plus 2pt minus 1pt}
\allowdisplaybreaks[2]
\newcommand{\E}{\mathbb E}
\newcommand{\R}{\mathbb R}
\newcommand{\C}{\mathbb C}
\newcommand{\ind}{\mathbf 1}
\newcommand{\heading}[1]{\par\smallskip{\large\bfseries\color{ink}#1}\par\smallskip}
\newcommand{\subhead}[1]{\par\smallskip{\bfseries\color{ink}#1}\enspace}
\newcommand{\qed}{\hfill\(\square\)}
\hypersetup{pdftitle={Gaussian domination in the charge-one cubic chaos}}
\begin{document}
\thispagestyle{plain}
{\LARGE\bfseries\color{ink} Gaussian domination in the\\[2pt] charge-one cubic chaos\par}
\vspace{3pt}
{\color{muted} A counterexample for unrestricted parameters and two comparison theorems\par}
\medskip

Let \(T=|z_1|^2\), \(a=\ell-2\beta>0\), \(L=z_1(a+\beta T)\), and
\[
 \Lambda(G)=\E|L+G|+\E|L-G|-\varepsilon^{-1}\E|G|^2,\qquad
 S=\sup_{\sigma\ge0}\Lambda(\sigma W).
\]
All complex standard Gaussians have second absolute moment one. Write
\(B(v)=2\E|L+\sqrt v\,W|\) for \(v\ge0\), with \(W\) independent of the inputs.

\textbf{Results.} A specific \(G\in H_{2,1}(\C^2)\) and positive \((\varepsilon,\ell,\beta)\) give
\(\Lambda(G)>S\). In the application box we prove same-norm domination for conditionally
proper Gaussian perturbations with the variance profiles in Section 3. The family includes a
remote cubic that shares a Gaussian coordinate with \((T-1)z_2\).
Section 4 proves the scalar Stein-sign inequality at every scale for the real cubic \(H_3(X)\).
The uniform \(10^{-4}\) bound for arbitrary cubic tensors in the box remains open here.

\heading{1. Which variances can be Gaussian optimizers?}
The density of \(L\) on the plane is radial. At \(r(t)=\sqrt t\,(a+\beta t)\) its value is
\[
 p_L(r(t))=\frac{e^{-t}}{\pi(a+\beta t)(a+3\beta t)}.
\]
This follows from \(T\sim\mathrm{Exp}(1)\), the uniform phase, and the derivative of \(r(t)^2\).
In particular, \(p_L\) is strictly decreasing in the radius. The planar heat equation with
generator \(\Delta/4\), together with \(\Delta|x|=|x|^{-1}\) in distributions, gives
\[
 B'(v)=\tfrac12\E |L+\sqrt v\,W|^{-1}\quad(v>0),\qquad
 B'(0+)=\tfrac12\int_0^\infty\frac{e^{-t}}{\sqrt t(a+\beta t)}\,dt.
\]
The convolution \(p_L*|\cdot|^{-1}\) is continuous, finite and strictly radially decreasing.
For the monotonicity, write both factors by their disk superlevel sets. The area of overlap
of two disks decreases with the distance between their centers, and for any two distinct
distances there is a positive-measure set of radii for which it decreases strictly.
Thus \(B'(v)\) decreases strictly by coupling the radii \(\sqrt v\,|W|\); it tends to zero
as \(v\to\infty\), since \(\sup_x\E|x+\sqrt v W|^{-1}=\sqrt{\pi/v}\).
Consequently \(B\) is strictly concave. For each \(v>0\), the choice
\[
 \boxed{\quad\varepsilon=1/B'(v)=2\big/\E |L+\sqrt v\,W|^{-1}\quad}
\]
makes \(v\) the unique maximizing variance. For fixed \(a,\beta\), domination for
\emph{every} penalty is therefore equivalent to the same-norm comparison at every variance.

The application box includes a zero optimizer. At
\((\varepsilon,\ell,\beta)=(0.74,0.90,0.12)\),
\(B'(0+)<\sqrt\pi/(2\cdot0.66)<1/0.74\), so \(\sigma_*=0\).
For comparison, numerical quadrature of the derivative gives
\(\sigma_*\simeq0.35601144\) at \((0.76,0.76,0.16)\) and
\(\sigma_*\simeq0.20086017\) at \((0.7564,0.8456,0.1432)\).
Only the zero-optimizer example is used in this argument.

\newpage
\heading{2. An explicit two-coordinate counterexample}
Take \(\sigma=10^{-15}\), \(a=\sigma^2\), \(\beta=1\), \(\ell=2+a\), and choose
\(\varepsilon=1/B'(\sigma^2)\) for this profile. Put \(X=z_2,\ U=|X|^2,\ d=1/12\), and
\[
 P(t,X)=\frac{X[(t-1)+d(U-2)]}{\sqrt{1+2d^2}},\qquad G=\sigma P(T,X).
\]
The two summands are Wick cubics, orthogonal with squared norms \(1\) and \(2\)
before multiplication by their coefficients. Hence \(\|G\|_2=\sigma\).
Replacing \(X\) with \(-X\) gives the equality of the two signs in \(\Lambda(G)\).

\subhead{The moment at \(t=0\).}
Write \(q=5/3\) and \(M=\E|W|^q=\Gamma(1+q/2)\).
The tangent bound \(|1-dx|^q\ge1-qdx\) and the identity
\(\E U^{q/2}(U-2)=(q/2-1)M\) give
\[
 \frac{\E|P(0,X)|^q}{M}
 \ \ge\ \frac{1+5d/18}{(1+2d^2)^{5/6}}
 \ \ge\ \frac{1+5d/18}{1+5d^2/3}=\frac{442}{437}.
\]
Also \(\E|P(0,X)|^2=1\). Since \(M\ge\mathbb P(|W|\ge1)=1/e\),
the excess in the \(q\)-moment exceeds \(1/300\).

\subhead{Converting the moment excess into a norm excess.}
Let \(h(r,y)=\E_\theta|r e^{i\theta}+y|-r\), with \(\theta\) uniform and \(r,y\ge0\).
Then \(0\le h(r,y)\le\min(y,y^2/(2r))\) for \(r>0\), by Jensen and the triangle inequality.
Convexity in \(y\) also gives
\(0\le \partial_y^+h(r,y)\le[h(r,2y)-h(r,y)]/y\le2y/r\).
In particular \(|h(r,y)-h(r,w)|\le |y^2-w^2|/r\). By homogeneity and scaling,
\[
 \int_0^\infty h(t^{3/2},\sigma y)\,dt=C\,\sigma^{5/3} y^{5/3},\qquad
 C=\int_0^\infty h(t^{3/2},1)\,dt>1/5 .
\]
The integral is finite by the bounds above. For \(t\le1/4\), Jensen gives
\(h(t^{3/2},1)\ge 1-t^{3/2}\ge7/8\), proving the lower bound on \(C\).

First set \(a=0\), so \(|L|=T^{3/2}\). Averaging the phase gives the gain
\[
 \E|L+G|-\E|L|=\int_0^\infty e^{-t}\E h(t^{3/2},\sigma|P(t,X)|)\,dt .
\]
Replacing \(P(t,X)\) by \(P(0,X)\) costs at most
\[
 \sigma^2\int_0^\infty e^{-t}\frac{2t+t^2}{t^{3/2}}\,dt
 =\tfrac52\sqrt\pi\,\sigma^2 .
\]
Indeed \(P(t,X)-P(0,X)=tX/\sqrt{1+2d^2}\), and Cauchy--Schwarz bounds the expected
absolute difference of squared moduli by \(2t+t^2\). Replacing \(e^{-t}\) by \(1\)
costs at most \(\tfrac12\sigma^2\int_0^\infty(1-e^{-t})t^{-3/2}\,dt=\sqrt\pi\,\sigma^2\).
For the Gaussian comparator, using \(1\) in place of \(e^{-t}\) gives an upper bound.
Finally restoring \(a=\sigma^2\) changes the difference of the two mean norms by at
most \(2a\E|z_1|=\sqrt\pi\,\sigma^2\). Altogether,
\[
 \E|L+G|-\E|L+\sigma W|
 > \frac{\sigma^{5/3}}{1500}-9\sigma^2>0 .
\]
The equal penalties and Section 1 now give \(\Lambda(G)>\Lambda(\sigma W)=S\).
The witness has \(\ell-2\beta=10^{-30}\), outside the required box.\qed

\newpage
\heading{3. A same-norm theorem in the application range}
Suppose \(0<a/\beta\le11/2\). This contains the entire stated box.
Let \(\mathcal F\) be independent of \(z_1\). Assume that, conditionally on
\((z_1,\mathcal F)\), \(G\) has a centered proper complex Gaussian law of variance
\[
 h(T)=s+u(T-1)^2+bT+cT^2,\qquad s,u,b,c\ge0.
\]
The coefficients can be arbitrary \(\mathcal F\)-measurable random variables with
\(\E(s+u+b+2c)<\infty\). The Gaussian can be degenerate. Then
\[
 \boxed{\ \E|L+G|+\E|L-G|\ \le\ 2\E|L+\|G\|_2 W|.\ }
\]
In particular \(\Lambda(G)\le S\) for every \(\varepsilon>0\).
It suffices to prove the claim for deterministic coefficients. Indeed, the conditional
mean variance is \(\bar h=s+u+b+2c\); conditional comparison followed by Jensen
for the concave \(B\) gives \(\E B(\bar h)\le B(\E\bar h)\).

\subhead{Two bounds on the inverse-radius kernel.}
For \(w>0\), let \(M_w(r)=\E|r+\sqrt w W|\) and \(K_w(r)=\E|r+\sqrt w W|^{-1}\).
The heat equation gives \(\partial_w M_w=K_w/4\). With \(x=r^2/(2w)\),
\[
 K_w(r)=\sqrt{\pi/w}\,e^{-x}I_0(x),\qquad
 m(x)=2x\big(1-I_1(x)/I_0(x)\big).
\]
The kernel formula follows by integrating the Gaussian Laplace transform against
\(y^{-1/2}dy/\sqrt\pi\) and substituting \(wy/(1+wy)\). Logarithmic differentiation gives
\[
 \partial_{\log r}\log K=-m,\qquad
 \partial_{\log w}\log K=(m-1)/2,\qquad
 m'=2-2m+m^2/(2x).
\]
We use the elementary estimates
\[
 0<m(x)<13/10,\qquad \sup_{r,w>0} r K_w(r) <6/5 .
\]
For the first one, put \(D=13/10,\ x_*=D^2/[4(D-1)]=169/120\).
The recurrence \(I_0-I_2=2I_1/x\) and the series estimate \(I_2/I_1<x/4\) give
\(I_1/I_0>x/(2+x^2/4)\). At \(x=x_*\) this lower bound exceeds
\(7/13=1-D/(2x_*)\), so \(m(x_*)<D\).
At a crossing of \(D\), the differential equation gives \(m'>0\) to the left
of \(x_*\) and \(m'<0\) to its right. Since \(m(0)=0\), \(m\) stays below \(D\).
Positivity follows from \(I_1<I_0\).

For the second estimate, \(rK_w(r)=\sqrt{2\pi x}\,e^{-x}I_0(x)\), whose logarithmic
derivative in \(x\) is \((1-m)/(2x)\). The series gives \(I_1(1)<I_0(1)/2\), hence
\(m(1)>1\). Every crossing of \(1\) has positive derivative, so the ratio decreases
for \(x\ge1\). On \(0<x\le1\), use \(I_0(x)\le e^{x^2/4}\).
The upper bound \(\sqrt{2\pi x}\,e^{-x+x^2/4}\) is increasing there, with value
\(\sqrt{2\pi}e^{-3/4}<6/5\) at \(x=1\).
The same crossing argument shows that as a function of \(w\), \(K_w(r)\) has at most
one interior extremum, a maximum. Hence for \(0\le w\le w_0\) and \(r>0\),
\[
 K_w(r)\ \ge\ \min\{1/r,K_{w_0}(r)\}.
\]
Here \(K_0(r)=1/r\) by the zero-variance limit. The Bessel identities used above are
also recorded in NIST DLMF, \href{https://dlmf.nist.gov/10.25}{\S10.25},
\href{https://dlmf.nist.gov/10.29}{\S10.29}, and \href{https://dlmf.nist.gov/10.32}{\S10.32}.

\newpage
\heading{Proof of the comparison}
First replace the constant variance contribution \(u\) by \(u(t-1)^2\) along the path
\[
 w_d(t)=s+bt+ct^2+u\{1-d+d(t-1)^2\},\qquad 0\le d\le1.
\]
The derivative of the mean norm is
\[
 \frac u4\int_0^\infty e^{-t}t(t-2)K_{w_d(t)}(r(t))\,dt .
\]
For \(u>0\), differentiation is justified on each compact subinterval of \(d<1\)
by \(w_d\ge u(1-d)\). Continuity gives the comparison at \(d=1\).
The integral of \(e^{-t}t(t-2)\) is zero. We show the kernel is at least its value
at \(t=2\) for \(t<2\), and at most that value for \(t>2\).

For \(t\ge3/2\), \(w_d'(t)\ge0\) and \(w_d'/w_d\le 2/(t-1)\). This holds term by
term for the nonnegative summands. Also,
\[
 (t-1)\frac{r'(t)}{r(t)}
 =\frac{t-1}{2t}+\frac{t-1}{a/\beta+t}
 \ \ge\ \frac16+\frac1{14}=\frac5{21}>\frac3{13}.
\]
Therefore \(K_{w_d(t)}(r(t))\) decreases in \(t\) on this interval. In detail, its
logarithmic derivative is \(-m r'/r+(m-1)w_d'/(2w_d)\). It is negative when \(m\le1\).
When \(1<m<13/10\), use \((m-1)/m<3/13\) and the two preceding estimates.

For \(0<t\le3/2\), we have \(w_d(t)\le w_d(2)\) and
\[
 \frac{r(t)}{r(2)}\le\frac{r(3/2)}{r(2)}
 \le\frac{7\sqrt3}{15}<\frac56 .
\]
At \(w=w_d(2)\) the kernel decreases with \(r\). The bounds on the preceding page
also give \(1/r(t) \ge (6/5)/r(2)>K_{w_d(2)}(r(2))\). Taking the minimum of these
two endpoint bounds proves the desired comparison. It follows that the derivative
in \(d\) is nonpositive. For \(G=\alpha(T-1)z_2\) alone, taking \(d=\rho^2\) proves
monotonicity along the stated Ornstein--Uhlenbeck interpolation at every norm.

We are left with the variance \(s+u+bt+ct^2\). Interpolate this from its mean
\(\bar h\) by the affine path
\(\widetilde w_d(t)=(1-d)\bar h+d(s+u+bt+ct^2)\).
Here \(0\le\widetilde w_d'/\widetilde w_d\le2/t\), and \(r'/r\ge1/(2t)\).
The logarithmic derivative of the kernel is again negative: for \(m>1\) it is
at most \((m/2-1)/t<0\), and for \(m\le1\) its sign is immediate.
The derivative of the mean norm is therefore nonpositive, since
\(b(t-1)+c(t^2-2)\) is increasing with exponential mean zero, and the kernel is
decreasing. Both variance replacements decrease the mean norm when moving away
from the constant variance. This proves the theorem.\qed

\subhead{A shared remote coordinate.}
For any real \(\alpha,\lambda\), consider
\[
 G=z_2\{\alpha(|z_1|^2-1)+i\lambda(|z_3|^2-1)\}.
\]
This is a sum of two charge-one Wick cubics. The remote cubic and the visible-dependent
cubic share \(z_2\). Given \(z_1,z_3\), their sum is proper Gaussian of variance
\(\alpha^2(T-1)^2+\lambda^2(|z_3|^2-1)^2\), so the theorem applies at every amplitude.
The example with \(\lambda=0\) already gives a nonconstant, nonmonotone conditional
variance. For a general cubic the conditional Gaussian hypothesis can fail, for example
when a remote term is \(z_2(|z_2|^2-2)\) and shares \(z_2\) with \((T-1)z_2\).

\newpage
\heading{4. The real cubic \(H_3\) at every scale}
Let \(X\sim N(0,1)\), \(Y=H_3(X)=X^3-3X\), and
\(A=|DY|^2/3=3(X^2-1)^2\). Then \(\E A=\E Y^2=6\). In fact,
\[
 \boxed{\qquad \E[(A-6)\ind_{\{|Y|\le h\}}]<0\quad\text{for every }0<h<\infty .\qquad}
\]
By layer cake this gives \(\operatorname{Cov}(A,\varphi(bY))<0\) for every \(b\ne0\).
Since conditioning \(A\) on \(Y\) preserves this covariance, it gives precisely the
Stein-kernel inequality at every scale for every nonzero element of the third chaos
on one real Gaussian coordinate.

To prove the boxed statement, put \(F(x)=x(x^2+1)e^{-x^2/2}\). Up to a positive
constant, the signed density on the positive half-line is
\[
 (x^4-2x^2-1)e^{-x^2/2}=-F'(x).
\]
If \(h\ge2\), the positive sublevel set is one interval \([0,r]\); its signed mass
is \(-F(r)<0\). For \(0<h<2\), it is \([0,r_1]\cup[r_2,r_3]\), with
\[
 0<r_1<1<r_2<\sqrt3<r_3<2,\qquad r_3=r_1+r_2.
\]
Indeed \(r_1,r_2,-r_3\) are the three roots of \(x^3-3x+h=0\).
The signed mass is therefore \(-F(r_1)+F(r_2)-F(r_3)\).

On \([0,1]\), \((1+x^2)e^{-x^2/2}\ge1\), so \(F(r_1)\ge r_1\).
On \([1,2]\), \(-F'(x)\le7e^{-2}<1\).
For this bound, differentiate \(-F'\); the derivative has the sign of
\(x(-x^4+6x^2-3)\), which is positive on this interval. Consequently
\[
 F(r_2)-F(r_3)\le7e^{-2}(r_3-r_2)<r_1\le F(r_1).
\]
Symmetry gives the claim on the full line.\qed

\subhead{Relation to the formulation with \(\Phi\).}
For \(Y_b=bY\), Gaussian integration by parts in the third chaos gives
\(\E[Y_b f(Y_b)]=\E[b^2 A f'(Y_b)]\). Applying this to \(f(y)=2\Phi(y)-1\) yields
\[
 \E\!\left[|Y_b|(2\Phi(|Y_b|)-1)\right]
 < 2\,\E[Y_b^2]\,\E[\varphi(Y_b)]\qquad(b\ne0).
\]
The rank-one result allows a direct proof even though the cubic has turning points
at \(x=\pm1\). A different family in any dimension also satisfies the inequality:
if \(Y\) is conditionally centered real Gaussian with random variance \(V\), integrating
by parts conditionally gives
\[
 \begin{aligned}
 &\E\!\left[|Y|(2\Phi(|Y|)-1)\right]-2\,\E[Y^2]\,\E[\varphi(Y)]\\
 &\hspace{18mm}=\frac{2}{\sqrt{2\pi}}\operatorname{Cov}
       \left(V,(1+V)^{-1/2}\right)\le0 .
 \end{aligned}
\]
This covers, for example, cubics linear in an independent block of real Gaussian
coordinates, with quadratic coefficients in the other block. The scalar inequality
for general third-order tensors and general higher orders remains unresolved here.

For the original complex problem at \(L=0\), \(S=\pi\varepsilon/4\).
Optimizing \(\Lambda(bG)\) over real \(b\ge0\) shows that this endpoint of K is exactly
the assertion \(\E|G|\le(\sqrt\pi/2)\|G\|_2\) for the charge-one cubic chaos.
The all-scale scalar Stein-sign conjecture would additionally control the interpolation.
\end{document}
